% ===================================================================
%  TAULUKKO N:o 11 — Kaupunki ja sen muistomerkit. — Tulipalo.
%  Traduction 2026 d'apres texto/io, controlee sur texto/fr et
%  texto/sv. Consigne : voir texto/fi/10-taulukko-01.tex.
%
%  TROIS INVERSIONS, la meme forme que partout : « la pordo (8) dil
%  olda fortifikita kastelo (7) », « la kupolo (40) dil hospitalo
%  (39) », « la muzikisti (81) dil orkestro (80) ». Relative dans les
%  trois cas. La premiere n'a PAS ete vue a la redaction : c'est le
%  garde qui l'a trouvee, et c'est a cela qu'il sert.
%
%  ET UNE TROISIEME QUI NE S'EST PAS PRESENTEE : au bloc c2-02-1,
%  l'ido met la comparaison avant l'objet — « plu bone kam la soldati
%  (57)... la kuriozi (56) » — ce qui avait coute une refonte au
%  neerlandais. Le finnois range sa phrase par le theme et non par la
%  fonction : l'ordre de l'ido s'y ecrit tel quel.
% ===================================================================

%%K t11-apar-1 apar
\VUsaut{2.0mm}
\VUcentre{13.2pt}{90}{KOLMAS SARJA}
\VUsaut{3.0mm}
\VUfilet{16mm}
\VUsaut{5.6mm}
\VUtitre{15.0pt}{60}{TAULUKKO N:o 11}
\VUsaut{1.4mm}
\VUpk{11.4pt}{40}{Kaupunki ja sen muistomerkit. --- Tulipalo.}
\VUsaut{2.2mm}

%%K t11-tit-1 sub
\VUpk{10.2pt}{40}{Esikaupunki.}
\VUsaut{3.0mm}

%%K t11-c1-01-1 p
1. --- Asun suuressa kaupungissa, joka on 100 kilometrin päässä valtamerestä ja
joka kuitenkin on ensiluokkainen \VUgras{satama} \textsuperscript{(1)}. Joki,
joka sitä rajaa, on 400 metriä leveä ja muodostaa hyvin kauniin redin.

%%K t11-c1-02-1 p
2. --- Koko se kaupungin osa, joka on \VUgras{laitureiden}
\textsuperscript{(2)} varrella, näyttää aivan merenkulun- ja
teollisuudenvoittoiselta ja muodostaa \VUgras{esikaupungin}
\textsuperscript{(3)}, jota asuttavat vain merimiehet ja työläiset. Nämä
työskentelevät \VUgras{tehtaissa} \textsuperscript{(4)} ja
\VUgras{kaasulaitoksessa} \textsuperscript{(5)}, jonka korkea \VUgras{savupiippu}
\textsuperscript{(6)} ja \VUgras{kaasukello} näkyvät hyvin kauas.

%%K t11-c1-03-1 p
3. --- Keskiajalla kaupunki oli linnoitettu, ja sen valleista löytää vielä
muutamia jäänteitä, varsinkin \VUgras{portin} \textsuperscript{(8)}, joka
johtaa vanhaan \VUgras{linnoitettuun linnaan} \textsuperscript{(7)} ja jonka
sivustoilla ovat sen kaksi \VUgras{tornia} \textsuperscript{(9)}, jotka nyt
toimivat \VUgras{vankilana}.

%%K t11-c1-04-1 p
4. --- Koko siinä \VUgras{kaupunginosassa}, joka näyttää hyvin vanhalta, vain
yksi rakennus on suhteellisen nykyaikainen : se on \VUgras{tullitalo}
\textsuperscript{(10)}. Se on laiturin ja sen \VUgras{pienen kadun}
\textsuperscript{(11)} välissä, joka johtaa vanhalle \VUgras{raatihuoneelle}
\textsuperscript{(12)}. Tuon viimeksi mainitun muistomerkin tuntee helposti sen
jättiläismäisestä \VUgras{kellotornista} \textsuperscript{(13)}, joka hallitsee
vanhaa kaupunkia.

%%K t11-c1-05-1 p
5. --- Kadun toisella puolella \VUgras{seisoo} rakennus, joka on pystytetty
samaan tyyliin kuin tullitalo : se on \VUgras{Pörssi} \textsuperscript{(14)}.
Yksi sen julkisivuista antaa pienelle torille, jonka varrella
\VUgras{lääninhallitus} \textsuperscript{(15)} on. Kaupunkimme on nimittäin,
olematta pääkaupunki, kuitenkin tärkeä läänin pääkaupunki.

%%K t11-tit-2 sub
\VUsaut{5.0mm}
\VUpk{10.2pt}{40}{Kaupungin korkein osa.}
\VUsaut{3.4mm}

%%K t11-c1-06-1 p
6. --- Varuskunta, jokseenkin lukuisa, on majoitettu laajoihin hyvin terveellisiin
\VUgras{kasarmeihin} \textsuperscript{(16)}; ne ulottuvat aina
\VUgras{kaupunginpuiston} \textsuperscript{(17)} ristikkoaitaan asti.
\VUgras{Bulevardi} \textsuperscript{(19)}, joka johtaa tuohon puistoon, kulkee
\VUgras{säästöpankin} \textsuperscript{(18)} takaa ja päättyy
\VUgras{tuomiokirkon} \textsuperscript{(20)} torille.

%%K t11-c1-07-1 p
7. --- Kaikki nuo muistomerkit ulottuvat amfiteatterin tavoin kukkulalla, jolla
myös laaja \VUgras{oppilaitoksemme} \textsuperscript{(21)} on.

%%K t11-c1-07-2 p
Jos laskeudutaan hieman viettävää tietä pitkin, joka yhtyy Suurkatuun, tullaan
\VUgras{Oikeustaloon} \textsuperscript{(22)}, jonka edessä ulottuu miellyttävä
\VUgras{julkinen puutarha} \textsuperscript{(26)}. Tuo \VUgras{istutus} ei ole
kovin suuri, mutta se muodostaa keskelle kaupunkia vehreyden ja viileyden
keitaan. Siksi siellä käyvät paljon kaupunkilaiset, jotka tulevat mielellään
istumaan \VUgras{kahvilan} \textsuperscript{(25)} telttakatoksen alle
kuuntelemaan sotilassoittajia, jotka torstaisin ja sunnuntaisin soittavat
\VUgras{kioskissa} \textsuperscript{(27)}, joka on nurmikoiden ja pensaikkojen
välissä.

%%K t11-c1-08-1 p
8. --- Yhdellä noista nurmikoista on pronssinen \VUgras{patsas}
\textsuperscript{(28)}, muistomerkki, joka on pystytetty erään kaupunkilaisemme
muistoksi, joka teki suuria palveluksia syntymäkaupungilleen.

%%K t11-tit-3 sub
\VUsaut{4.0mm}
\VUpk{10.2pt}{40}{Kulkuneuvot.}
\VUsaut{3.4mm}

%%K t11-c1-09-1 p
9. --- Tähän asti kaupunkiamme ei vielä valaista sähkövalolla, siinä on vain
\VUgras{kaasulyhtyjä} \textsuperscript{(29)}. Mutta \VUgras{paikallislehdet}
käyvät kampanjaa sen puolesta, että tuota valaistusjärjestelmää parannettaisiin,
\VUgras{samoin kuin} on jo täydellistetty \VUgras{raitiovaunujen
vetojärjestelmä}. \VUgras{Ne} vanhat vaunut, joita \VUgras{hevoset}
\VUgras{vetivät}, on käytännössä korvattu \VUgras{sähköraitiovaunuilla}
\textsuperscript{(31)}, jotka nopeasti vievät matkustajat \VUgras{kaupungin}
päästä \VUgras{toiseen} hyvin \VUgras{halvalla} hinnalla.

%%K t11-c1-10-1 p
10. --- Oikeustalon luona on \VUgras{Pankki} \textsuperscript{(32)}, jossa
diskontataan kauppapaperit. \VUgras{Pankkivirkailijat} \textsuperscript{(33)}
käyvät perimässä velat kodeissa.

%%K t11-tit-4 sub
\VUsaut{4.0mm}
\VUpk{10.2pt}{40}{Taide ja opetus.}
\VUsaut{3.4mm}

%%K t11-c1-11-1 p
11. --- Kaunis rakennus, joka on aivan vieressä, on \VUgras{museo}. Se sisältää
samalla kertaa kaupungin \VUgras{kirjaston} \textsuperscript{(34)}, kauniit
\VUgras{luonnontieteelliset kokoelmat} \textsuperscript{(35)}
(Luonnonhistoriallisen museon) ja sievän \VUgras{taulumuseon}
\textsuperscript{(36)}, joka muodostaa komean pysyvän \VUgras{näyttelyn}.

%%K t11-c1-11-2 p
Se avataan yleisölle kahdesti viikossa, mutta vieraat saavat käydä siellä minä
päivänä hyvänsä \VUgras{vahtimestarille} \textsuperscript{(37)} annettua
juomarahaa vastaan. \VUgras{Maalarit} \textsuperscript{(38)} tulevat sinne
työskentelemään. Heidät näkee maalaustelineensä \VUgras{edessä},
\VUgras{paletti} kädessään, jäljentämässä kuuluisia tauluja.

%%K t11-c1-12-1 p
12. --- Kukkulalla, museon takana, alamme nähdä \VUgras{kupolin}
\textsuperscript{(40)}, joka on \VUgras{sairaalan} \textsuperscript{(39)}
päällä, ja \VUgras{seminaarin} \textsuperscript{(41)} rakennukset, jossa
kunnallisten koulujemme opettajat koulutettiin. Teatteritorilla on
\VUgras{postikonttori} \textsuperscript{(43)}, joka on sijoitettu
\VUgras{yksityistaloon} \textsuperscript{(42)}.

%%K t11-tit-5 sub
\VUsaut{5.0mm}
\VUpk{11.4pt}{40}{Tulipalo.}
\VUsaut{3.0mm}

%%K t11-tit-6 sub
\VUpk{10.2pt}{40}{Palohälytys.}
\VUsaut{3.4mm}

%%K t11-c2-01-1 p
1. --- Hirvittävä huhu leviää kaupungin läpi : teatterissa on juuri syttynyt
tulipalo! Sillä hetkellä, kun tulen \VUgras{torille} \textsuperscript{(96)},
väkijoukko on jo koolla \VUgras{jalkakäytävällä} \textsuperscript{(46)} ja
\VUgras{ajoradalla} \textsuperscript{(47)}. \VUgras{Postivirkailijat}
\textsuperscript{(44)} ja \VUgras{kirjeenkantajat} \textsuperscript{(45)} tulevat
ulos postikonttorista. \VUgras{Raitiovaununkuljettajan} \textsuperscript{(48)}
täytyy pysäyttää (raitio)vaununsa antaakseen kunnallisen
\VUgras{ambulanssivaunun} \textsuperscript{(50)} kulkea ohi, joka tulee
\VUgras{kirurgin} \textsuperscript{(52)} ja \VUgras{sairaanhoitajan}
\textsuperscript{(51)} kanssa hoitamaan \VUgras{haavoittunutta}
\textsuperscript{(53)}, joka on kannettu \VUgras{paareilla}
\textsuperscript{(54)} \VUgras{sanomalehtikioskille} \textsuperscript{(55)}.

%%K t11-c2-02-1 p
2. --- Jalkaväen komppania, joka palasi harjoituksista, pysähtyi turvatakseen
järjestyksenpidon yhdessä \VUgras{ratsastavan kaupunginpoliisin}
\textsuperscript{(61)} kanssa. \VUgras{Ratsastajat}, joiden hevoset nousevat
takajaloilleen, saavat paremmin kuin \VUgras{sotilaat} \textsuperscript{(57)}
tai \VUgras{santarmit} \textsuperscript{(63)} \VUgras{uteliaat}
\textsuperscript{(56)} perääntymään. Muutoinkin santarmeilla on paljon tekemistä.
Yksi heistä näyttää \VUgras{poliisikonstaapeleille} \textsuperscript{(64)},
minne on vietävä toinen \VUgras{loukkaantunut} \textsuperscript{(65)}, urhea
palomies, joka on pudonnut tikkailta.

%%K t11-c2-03-1 p
3. --- \VUgras{Upseeri} \textsuperscript{(66)} osoittaa tarkoin sen paikan, johon
on asetettava \VUgras{höyryruisku} \textsuperscript{(67)}, joka on tulossa.
Odottaessaan tuota voimakasta konetta hän on pystyttänyt \VUgras{käsiruiskun}
\textsuperscript{(68)}, jonka pitkä \VUgras{letku} \textsuperscript{(69)}
heittää vettä kuin tulvavirtana tulelle. Kuusi palomiestä käyttää sitä, samalla
kun heidän toverinsa, joka pitää \VUgras{suihkuputkea} \textsuperscript{(70)},
suuntaa \VUgras{suihkun} \textsuperscript{(71)} palon keskipisteeseen.

%%K t11-c2-04-1 p
4. --- Teatterin toisella puolella on pystytetty suuret \VUgras{tikkaat}
\textsuperscript{(72)}. Ne tekevät palomiehille mahdolliseksi lähestyä liekkejä,
jotka lyövät ylös mustan \VUgras{savun} \textsuperscript{(75)} pyörteiden
välistä.

%%K t11-tit-7 sub
\VUsaut{5.0mm}
\VUpk{10.2pt}{40}{Urheita tekoja.}
\VUsaut{3.4mm}

%%K t11-c2-05-1 p
5. --- \VUgras{Tulipalo} \textsuperscript{(74)} syttyi \VUgras{teatterin}
\textsuperscript{(73)} lämpiössä, samalla kun harjoiteltiin \VUgras{oopperaa},
jonka ensimmäinen \VUgras{näytäntö} olisi pitänyt olla huomenna. Onneksi
salongissa oli vain vähän \VUgras{katsojia} \textsuperscript{(76)}. Poloiset
ovat paenneet \VUgras{parvekkeelle} \textsuperscript{(77)}, jonne urheat
\VUgras{palomiehet} \textsuperscript{(86)} tulevat auttamaan heitä tikkailla ja
köysillä.

%%K t11-c2-06-1 p
6. --- \VUgras{Pylväskäytävän} \textsuperscript{(78)} alla, jossa
\VUgras{juliste} \textsuperscript{(79)} kertoo näytännöstä, näkee
\VUgras{soittajien} \textsuperscript{(81)}, jotka ovat \VUgras{orkesterissa}
\textsuperscript{(80)}, pakenevan kantaen soittimiaan : \VUgras{viulu}
\textsuperscript{(81)}, \VUgras{huilu} \textsuperscript{(82)},
\VUgras{sello} \textsuperscript{(83)}, \VUgras{pasuuna}
\textsuperscript{(84)}, \VUgras{rumpu} \textsuperscript{(85)}, j. n. e..
Koetetaan pelastaa osa \VUgras{kalustosta} \textsuperscript{(87)}.

%%K t11-c2-07-1 p
7. --- \VUgras{Näyttelijät} \textsuperscript{(89)} ja
\VUgras{näyttelijättäret} \textsuperscript{(88)}, \VUgras{laulajat} ja
\VUgras{laulajattaret} joutuvat pakenemaan \VUgras{teatteripuvussaan}
\textsuperscript{(90)}. Tenori selittää \VUgras{sanomalehtimiehelle}
\textsuperscript{(92)}, kuinka se tapahtui. \VUgras{Reportteri} riensi sinne
heti ensi hetkestä viranomaisten kanssa : \VUgras{maaherra}
\textsuperscript{(91)}, \VUgras{pormestari} \textsuperscript{(93)},
\VUgras{poliisimestari} \textsuperscript{(94)} ja \VUgras{oikeudellinen
virkamies} \textsuperscript{(95)}, jonka on tutkittava asia arvioidakseen
vastuun.
\VUsaut{6.0mm}
\VUfilet{22mm}
